\begin{table}[htbp]
\caption{Audit claim map: what the evidence supports, rejects, and fails to support.}
\label{tab:claim_status}
\centering
\small
\begin{tabular}{p{7.6cm}p{5.6cm}}
\toprule
Claim & Status \\
\midrule
Benign-only QLoRA training is valid one-class training & Supported \\
Fold-aligned detector strength reflects behavioral discrimination & Rejected (seen-vs-unseen-user effect) \\
r6.2 contains a sparse feature family active on profile tokens whose edits pass the audited estimands & Supported descriptively (best-candidate patching and necessity); held-out replication concentrates on the dominant user; necessity significant against proxy-matched controls (methods appendix) \\
r4.2 contains a sparse feature family active mainly on behavior tokens whose edits pass patching and ablation & Partly: the best-candidate patch contrast is positive (significant against proxy-matched controls), but on held-out users the mean over donors is negative; ablation replication persists under a benign-only dictionary and across 3B adapter seeds, patch-repair replication does not \\
Editing the selected r4.2 features changes behavior-token predictions more than editing controls & Supported for the joint edit under the SAE replacement procedure at strengths 0.75--1 on held-out users; not shown to exceed an equal-size perturbation at the same positions, or to be specific to anomalous days \\
The selected r4.2 edits act on behavior but not profile predictions & Not established: the profile effect is uncertain and the preregistered direct comparison is inconclusive \\
Editing features toward a benign colleague repairs malicious days & Not established: only the searched best candidate is positive; the mean over donors has the opposite sign \\
Configuration-independent r4.2 confirmation & Not established \\
Literal feature transfer across benchmarks succeeds & Rejected \\
Transfer failure is explained by SAE seed non-identifiability & Rejected (alignment controls) \\
Positive-population size alone explains the feature-selection dissociation (conditional on the fixed r4.2 dictionary) & Rejected at the selection stage (subsampling) \\
The profile/behavior attribution dissociation requires positive-shaped dictionaries & Rejected: benign-only SAE retrain reproduces it (r4.2 top-5 $\geq$94\% behavioral; r6.2 $\geq$99\% profile-bound in all LOUO folds) \\
The dissociation is particular to the 8B run, the serializer defect, or one adapter seed & Rejected: 3B repaired-serializer rerun reproduces attribution across two adapter seeds per benchmark, and r4.2 ablation dependence replicates across 3B seeds \\
The behavioral pole is an artifact of synthetic data or simulated psychometrics & Rejected: TWOS (real users, real Big-Five) yields ~100\% behavioral attribution across two adapter seeds at 50\% malicious prevalence, as the population account predicts \\
Profile capture is a generic dataset property that any detector class inherits & Rejected at attribution level: classical one-class detectors on the same r6.2 features place 4--7\% importance on profile, select no profile feature in any top-5, and lose nothing when profile is removed \\
The shortcut is an artifact of CERT's injected profile & Not tested cleanly (corrected 2026-09-24): on LANL, user sampling and fold assignment shared a hash, so the unseen pool is 23 red-team users with no negative-only user and the seen pool adds 988 negative-only users; the AUCs ($0.95$ seen, $0.50$ unseen; host-anonymized $0.88$, $0.63$) describe those pools only; a corrected split is pending \\
Profile content is merely correlated with the unseen-user collapse & Rejected in the audited 8B run: removing profile content from adaptation and scoring restores fold-aligned unseen-user ROC $0.53\to0.90$ (descriptive paired bootstrap interval excluding zero); small and non-significant in the audited 3B run (one adapter seed per arm) \\
\bottomrule
\end{tabular}
\end{table}
